\documentclass[a4paper,12pt]{article}
\usepackage[utf8]{inputenc} 

%===============================================================================================================================================
%Das gehört in die Preambel kopiert:
\usepackage[ngerman]{babel}
\usepackage{csquotes} % korrekte Anführungszeichen, wichtig für biblatex

\usepackage[
  backend=biber,      
  style=authoryear,                                             % Alternativen: numeric, ieee, apa, verbose, etc.
  giveninits=true,                                              % Initialen statt voller Vornamen
  uniquename=init,
  maxcitenames=2,                                               % „Müller & Meier“; danach „u. a.“
  maxbibnames=99,                                               % im Verzeichnis alles zeigen
  url=true, doi=true, isbn=false
]{biblatex}

\addbibresource{../Bücher/refs.bib}                                                                 %Dateipfad für die .bib-Datei
\DefineBibliographyStrings{ngerman}{andothers = {u.\,a.}, page = {S.}, pages = {S.}}
\ExecuteBibliographyOptions{sorting=nty}                                                            %Sortierung: Name, Titel, Jahr
%===============================================================================================================================================

\title{Literaturverzeichnis, Zitieren und BibTeX}
\author{SCHD, STOL, WACJ}
\date{13.10.2025}

\begin{document}
  \maketitle

  %Beispielzitat, weil im Literaturverzeichnis sonst nichts steht:
  \noindent
  Methodisch wurde das mehrfach diskutiert \parencite{einstein1905,knuth1984,lamport1994}.

  \printbibliography %Druckt das Literaturverzeichnis am Ende des Dokuments

\end{document}

%===============================================================================================================================================
%Zitiermöglichekeiten in LaTeX

%Autor im Satz:
\textcite{einstein1905} entwickelt die Spezielle Relativitätstheorie.

%Klammerzitat:
Klammerzitat: Das \LaTeX{}-System ist weit verbreitet \parencite{lamport1994}.


%Seiten- und Stellenangaben:
Eine Seite: \parencite[23]{knuth1984}. \\
Bereich: \parencite[23--27]{knuth1984}. \\
Kapitel/Abbildung: \parencite[Kap.\,2; Abb.\,3]{lamport1994}.

%Meherere Quellen auf einmal:
Methodisch wurde das mehrfach diskutiert \parencite{einstein1905,knuth1984,lamport1994}.

%Fußnotenzitate:
Webressourcen lassen sich bequem als Fußnote setzen:\footnote{Vgl.\ \footcite{w3cHTML}.}

%Autor oder Jahr einzeln:
Nur Autor: \citeauthor{einstein1905}; nur Jahr: \citeyear{einstein1905}.
%================================================================================================================================================

%================================================================================================================================================
%Literaturverzeichnis anpassen 

%Eigene Überschrift für Literaturverzeichnis, Eintrag ins Inhaltsverzeichnis:
\printbibliography[title={Literatur}, heading=bibintoc]

%Teil-Literaturverzeichnis nach Typ:
\printbibliography[type=book,     title={Bücher}]
\printbibliography[type=article,  title={Zeitschriftenartikel}]
\printbibliography[type=online,   title={Webquellen}]

%Teil-Literaturverzeichnis nach Stichwort:
\printbibliography[keyword=seminar, title={Quellen für das Seminar}]

% Familienname zuerst anzeigen (in Preambel):
\DeclareNameAlias{default}{family-given}

%Eigene Kategorien (z.b: Klassiker):
\DeclareBibliographyCategory{klassiker}
\addtocategory{klassiker}{knuth1984,lamport1994}
%================================================================================================================================================